% !TEX root = ../main.tex
%% Devoir surveillé n° 6 — chapitres 17 à 21 (valeurs vérifiées : verifications/devoirs.py)

\devoirDS{6}{Chapitres 17 à 21}{{La célérité des ultrasons : ondes et incertitudes}/7,{L'objectif d'un appareil photo : lentille mince}/7,{L'écran d'un smartphone : couleurs et photons}/6}

\begin{exoBac}{7}{La célérité des ultrasons}
Pour mesurer la célérité des ultrasons dans l'air, on utilise un émetteur de fréquence $f=\qty{40,0}{kHz}$ et deux récepteurs reliés à
un oscilloscope.
\begin{docu}{1}{Le protocole et les mesures}
Les deux récepteurs sont placés côte à côte : leurs signaux sont en phase. On éloigne lentement le second récepteur et l'on compte
dix nouvelles coïncidences des signaux : le récepteur s'est alors déplacé de dix longueurs d'onde, sur une distance $d$. Huit
groupes ont obtenu, en cm : 8,52 ; 8,48 ; 8,60 ; 8,45 ; 8,55 ; 8,50 ; 8,58 ; 8,47.
\end{docu}
\begin{docu}{2}{Influence de la température}
La célérité du son dans l'air dépend de la température : $v\approx331+0{,}6\,\theta$, en \unit{m.s^{-1}}, avec $\theta$ en \unit{\degreeCelsius}. Le jour de la
mesure, la salle était à \qty{15}{\degreeCelsius}.
\end{docu}
\begin{questions}
\item Les ultrasons sont-ils des ondes mécaniques ? Pourquoi ne les entend-on pas ?
\item Pourquoi mesure-t-on dix longueurs d'onde plutôt qu'une seule ?
\item Calculer la moyenne et l'écart-type de la série de mesures de $d$. En déduire la longueur d'onde.
\item Calculer la célérité des ultrasons.
\item Évaluer l'incertitude-type sur $d$ (type A), puis l'incertitude-type sur la célérité, et écrire le résultat.
\item \initiative Comparer qualitativement le résultat à la valeur attendue dans la salle, puis à la valeur souvent citée de
\qty{343}{m.s^{-1}}, valable à \qty{20}{\degreeCelsius}.
\end{questions}
\end{exoBac}

\begin{exoBac}{7}{L'objectif d'un appareil photo}
L'objectif d'un appareil photo est modélisé par une lentille mince convergente de distance focale $f'=\qty{50}{mm}$. Le capteur, de
\qty{24}{mm} sur \qty{36}{mm}, est placé à la distance de l'objectif où se forme l'image nette.
\begin{questions}
\item Calculer la vergence de l'objectif.
\item On photographie une personne de \qty{2,00}{m}, debout à \qty{5,00}{m} devant l'objectif. Calculer la position $\overline{\mathrm{OA'}}$ de son image.
\item Calculer le grandissement et la taille de l'image. La personne tient-elle entière, en hauteur, sur le capteur placé en
largeur ($\qty{24}{mm}$ de haut) ?
\item Pour photographier un paysage très éloigné, l'image se forme dans le plan focal image. De combien faut-il déplacer
l'objectif, et dans quel sens, pour passer de la personne au paysage ?
\item Sur le schéma ci-dessous, à l'échelle 1/2, construire l'image A'B' de l'objet AB. Mesurer $\overline{\mathrm{OA'}}$ et $\overline{\mathrm{A'B'}}$, puis vérifier
ces valeurs à l'aide de la relation de conjugaison.
\item Dans l'œil, quels éléments jouent le rôle de l'objectif et du capteur ?
\end{questions}
\schema{ds6-lentille}
\end{exoBac}

\begin{exoBac}{6}{L'écran d'un smartphone}
Chaque pixel de l'écran d'un smartphone comporte trois sous-pixels qui émettent de la lumière rouge (\qty{625}{nm}), verte (\qty{525}{nm})
et bleue (\qty{465}{nm}), d'intensités réglables.
\begin{donnees}
$h=\qty{6,63e-34}{J.s}$ ; $c=\qty{3,00e8}{m.s^{-1}}$ ; $\qty{1}{eV}=\qty{1,60e-19}{J}$.
\end{donnees}
\begin{questions}
\item L'écran utilise-t-il la synthèse additive ou la synthèse soustractive ? Justifier.
\item Quelle couleur perçoit-on lorsque seuls les sous-pixels rouge et vert sont allumés ? Lorsque les trois le sont ? Lorsque seuls
le vert et le bleu le sont ?
\item Calculer, en électronvolts, l'énergie d'un photon émis par chaque sous-pixel.
\item Calculer la fréquence et l'énergie, en joules, d'un photon vert.
\item Dans un sous-pixel, chaque photon est émis lors du passage d'un électron d'un niveau d'énergie à un niveau inférieur. Quel
sous-pixel met en jeu le plus grand écart entre niveaux ? Justifier.
\item Le soir, un « filtre de lumière bleue » diminue l'intensité des sous-pixels bleus. Quelle teinte prend alors une page blanche ?
\end{questions}
\end{exoBac}
